%!TEX root = ../main.tex

\chapter{Következtetések és jövőbeli irányok}

\section{Összegzés}

A dolgozatban a Zha rendszert mutattam be: egy hibrid szimbolikus zenegeneráló architektúrát, amely három különböző modellt oszt szét három különböző zenei aspektus között. A motiváció nem az volt, hogy egyetlen, mindent megoldó modellt építsek, hanem hogy a felelősségeket szétosszam ott, ahol a modellek erősségei természetes módon mások.

Három modul került integrálásra: a Markov-lánc (HMM-kiterjesztéssel) a hangnemtudatos monofón dallamért, a VAE (és a GOLC kiterjesztése) a transzpozíció-invariáns látens variációért, valamint a Transformer a polifón kíséretért, többsávos struktúráért és a szakaszhatárok közötti memória-átadásért. A három kimenet rögzített súlyozással keveredik, és skála- és regiszter-szűrésen áthaladva kerül MIDI fájlba.

\subsection{Hozzájárulások}

A dolgozat fő hozzájárulásai a következők:

\begin{enumerate}
  \item \textbf{Hibrid pipeline-szervezés}: a három modell explicit felelősség-szétosztása --- Markov a dallami alap, VAE a stílus-szintű variáció, Transformer a polifón kíséret --- egyetlen, REST API-n elérhető rendszerben.

  \item \textbf{GOLC-VAE módszer}: a Group-Orbital Latent Consistency regularizáció, amely a transzpozíció-invariáns látens reprezentáció tanulását egyetlen, könnyen implementálható veszteség-taggal éri el (lásd \ref{chap:vae}.\ fejezet); az implementáció a baseline VAE architektúráján belül marad, nem cseréli le a hálózati rétegeket csoport-ekvivariánsra.

  \item \textbf{Multitrack Transformer}: cross-attention koordináció a dallam-, basszus- és dobsávok között, három cross-attention modullal és track-típus embeddingekkel; a sávok harmóniai és ritmikai összehangolása így a hálózat belsejében történik, nem külső szabálykészleten.

  \item \textbf{Markov-modul kiterjesztései}: Laplace add-$\alpha$ simítás a ritka magasabb rendű táblákban, hőmérséklet- és top-K-vezérelt mintavételezés (egységesen a Transformer modullal), valamint határjelölő \texttt{START}/\texttt{END} szimbólumok a tanult kezdő és záró kontextusokhoz.

  \item \textbf{Zha Pipeline}: a három modell egymásra épülő integrációja, ahol a Markov-lánc dallami alapot ad, a VAE stílusvariációt biztosít, a Transformer pedig polifón kíséretet és hosszabb távú szerkezetet ad hozzá; a kimenetek keverése és szűrése egyetlen, végponttól végpontig futó pipeline-ban.
\end{enumerate}

\subsection{Gyakorlati eredmények}

A Lakh MIDI Dataset Clean adatbázis egy 17\,526 fájlból álló részhalmazán (megfelelő hosszúságú, $256 \leq$ tokens $\leq 4096$ MIDI fájlok) végeztem a tanításokat. A három modell külön-külön sikeresen tanult; a tanítási görbék (lásd \ref{fig:training_curves}.\ ábra) konvergálnak, és a generált kimenetek hangzásban koherensek a műfaji elvárásokkal. A baseline VAE és a GOLC-VAE összehasonlításában a \texttt{golc\_vae.py:222--226}-ban gyűjtött orbit-distance metrika közvetlen visszaesést mutat a $\beta_{\text{orbit}} > 0$ esetben, ami a transzpozíció-invariáns látens reprezentáció empirikus jele.

A teljes pipeline egységes, numerikus értékelését (egységesen mért perplexitás, harmonic coherence stb. minden modellkombinációra) a jelen dolgozat \emph{nem} közli, mert ennek reprodukálható futtatását időkereten belül nem stabilizáltam. A pipeline minőségi értékelése a kimentett MIDI fájlok (\texttt{output/} könyvtár) hallgatásával lehetséges, és a kísérleti munka során a kombinált pipeline subjective benyomása rendre jobb volt, mint bármelyik egyedi modellé --- különösen a basszus-dob koordinációban és a hangnem-konzisztenciában.

A generálási idő egy 30 másodperces darabra RTX 3090-en fél másodperc nagyságrendű; ez gyakorlatilag interaktív felhasználást enged.

\section{Korlátok és kihívások}

\subsection{Technikai korlátok}

\begin{enumerate}
  \item \textbf{VAE posterior-összeomlás}: a VAE a $\beta = 0{,}5$ konstans súlyozás mellett is alkalmanként a posterior-összeomlás közelébe kerül, ami a látens tér kifejezőképességét csökkenti. KL-annealing és free-bits regularizációval próbálkoztam, de a végleges konfigurációban ezek nem szerepelnek; a probléma teljes megoldása további munka.

  \item \textbf{Transformer memória- és számítási költség}: hosszabb szekvenciákon a self-attention $O(n^2 d)$ költsége és a memória puffer növekedése a generálási sebességet és a tárhelyet egyaránt megterheli, ami az interaktív felhasználás felső határát jelenti.

  \item \textbf{Rögzített súlyozás}: a három modell kimenetének $0{,}5 / 0{,}3 / 0{,}2$ keverési aránya empirikusan választott állandó érték. Egy bemeneti kontextustól függő, adaptív súlyozás valószínűleg jobb eredményt hozna; ennek megvalósítása jövőbeli irány.

  \item \textbf{Reprodukálható, egységes kiértékelés hiánya}: az egyes komponensekhez különálló metrika-szkriptek léteznek, de a teljes pipeline egyetlen, futtatható ablation-szkriptje nem készült el a jelen dolgozat időkeretén belül.
\end{enumerate}

\subsection{Zenei korlátok}

\begin{enumerate}
  \item \textbf{Műfaji eloszlás}: a Lakh MIDI Dataset Clean részhalmaza dominánsan pop és rock anyagot tartalmaz; a kevésbé reprezentált műfajokban (jazz, klasszikus, kortárs) a generált anyag stilisztikailag kevésbé meggyőző.

  \item \textbf{Hosszú formák}: a Transformer szekció-alapú memóriája segít a vers-refrén jellegű ismétlésekben, de a több perces zenei formák (szonáta-forma, rondó) konzisztens generálását nem oldja meg.

  \item \textbf{Dinamika és artikuláció}: a MIDI reprezentáció és a 8 csoportos velocity-felbontás óhatatlanul korlátozza az előadás-szintű kifejezőerőt.
\end{enumerate}

\section{Jövőbeli kutatási irányok}

A Zha jelen állapota több konkrét irányban továbbfejleszthető. Az alábbi felsorolás csak azokat az irányokat tartalmazza, amelyek a meglévő kódbázison közvetlenül megvalósíthatók.

\subsection{Rövid távú, közvetlenül következő munka}

\begin{enumerate}
  \item \textbf{Reprodukálható ablation-szkript}: egyetlen, parametrizálható kiértékelő, amely a teljes pipeline-on és minden részkombináción (Markov, VAE, GOLC-VAE, Transformer, illetve párosaikon) egységes metrikákat számol --- perplexitás, pitch-entrópia, harmonic coherence, és a $\mu$-vektorok orbit-távolsága. Ez tenné lehetővé a Zha számszerű összevetését más rendszerekkel.

  \item \textbf{Adaptív kimenet-súlyozás}: a Markov / VAE / Transformer kimenetek keverési arányának ($0{,}5 / 0{,}3 / 0{,}2$) bemeneti kontextustól függő, dinamikus paraméterezése. A jelenlegi rögzített súlyok egy egyszerű, a kontextus skála-egyértelműségétől függő szabályrendszerre cserélhetők.

  \item \textbf{Conditioning-bővítés}: a Transformer akkord- és tempó-conditioning mellé feltételes generálási opciók műfaj-, hangulat- vagy hangszerelés-jelölőkkel.

  \item \textbf{Attention-vizualizáció}: a Transformer attention súlyainak elemzése konkrét zenei részleteken; ez egyfelől interpretálhatósági munka, másfelől a head-redundancia diagnosztikája.
\end{enumerate}

\subsection{Hosszabb távú irányok}

\begin{enumerate}
  \item \textbf{Hierarchikus Transformer}: egy kétszintű architektúra, ahol az alsó szint a token-szintű, a felső szint a szakasz-szintű kontextust kezeli; ez a jelenlegi szekció-alapú memóriát váltaná le tanult mechanizmussal.

  \item \textbf{RLHF a sampling fölött}: az ismétlés-büntetés, top-K és top-p paraméterek hangolása emberi visszajelzés alapján; a Markov-modulban hasonló szabályok illesztését is lehetne RL-alapon.

  \item \textbf{Cross-modal kiterjesztés}: a szimbolikus generálás párosítása valamilyen audio-szintű (pl.\ \cite{copet2024simple} mintájára) komponenssel, hogy a kimenet ne csak MIDI, hanem renderelt hangzás is legyen.

  \item \textbf{Interaktív felület}: a jelenlegi FastAPI backend kiegészítése egy folyamatos, kotta-szintű interakcióra alkalmas frontenddel, ami valós időben fogadja a felhasználói módosításokat és újragenerálja az érintett szakaszt.
\end{enumerate}

\section{Záró megjegyzés}

A Zha rendszer egy mérnöki kísérlet arra, hogyan lehet több, eltérő természetű generatív modellt egyetlen, működő zenegeneráló pipeline-ba szervezni. A megközelítés nem azt állítja, hogy ez az egyetlen helyes felbontása a feladatnak; csak azt, hogy a felelősség-szétosztás --- skála-tudatos Markov a dallamért, transzpozíció-invariáns VAE a stílusért, többsávos Transformer a kíséretért --- egy konzisztens módon implementálható, és a kapott rendszer kimenetei zeneileg ellenőrizhetők. A jelen dolgozat ezt az implementációt írja le; a számszerű összevetés más rendszerekkel és a kiértékelés egységesítése a felsorolt rövid távú irányok között szerepel.
